% Point to the main file
\documentclass[../../Main.tex]{subfiles}

% ========== Start the document ==========
\begin{document}

\section{Section 1.4 - Propositional Functions and Quantifiers}

First, let's start with an example

\begin{example}
    Is $x + 3 = 10$ a proposition?
    
    \begin{solution}
        No, since $x$ is not know
    \end{solution}

    \medskip
    
    However, if $x$ is known, then $x + 3 = 10$ is a proposition

    The truth value depends on the value of $x$
\end{example}

\begin{definition}[Propositional Function]
    We call $x + 3 = 10$ a \textbf{propositional function} since the truth value of the proposition depends on the value of $x$
\end{definition}

\begin{example}[Examples of Propositional Functions]
    For the given propositonal functions, determine the truth value for the given value of variable(s)

    \begin{enumerate}

        \item 
            $P(x): x + 3 = 10$
    
            \begin{alphabet}
                \item
                {
                    $P(6)$

                    \begin{solution}
                        \FALSE since $6 + 3 \not= 10$
                    \end{solution}
                }
    
                \item
                {
                    $P(7)$

                    \begin{solution}
                        \TRUE since $7 + 3 = 10$
                    \end{solution}
                }
            \end{alphabet}

        \item 
            $P(a, b, c): a^2 + b^2 = c^2$

            \begin{alphabet}
                \item 
                    $P(11, 60, 61)$

                    \begin{solution}
                        \TRUE
                        \\
                        $11^2 + 60^2 = 121 + 3600 = 3721$
                        \\
                        $61^2 = (60 + 1)^2 = 3600 + 120 + 1 = 3721$
                    \end{solution}

                \item
                    $P(147, 201, 513)$

                    \begin{solution}
                        \FALSE
                        \\
                        $147^2 + 201^2 = \text{odd} + \text{odd} = \text{even}$
                        \\
                        $513^2 = \text{odd} \neq \text{even}$
                    \end{solution}
            \end{alphabet}

    \end{enumerate}
\end{example}

\begin{example}[Examples of Notation]
    \begin{enumerate}

        \item 
            Is "there exists a real number $x$ such that $x + 3 = 10$?" a proposition?
        
            \begin{solution}
                Yes, and this proposition is \TRUE since $x = 7$ works
            \end{solution}
        
            We can also rewrite this as:
        
            $$\exists x \in \mathbb{R} \mid x + 3 = 10$$

        \item 
            Is "for all real numbers $x$, $x + 3 = 10$" a proposition?

            \begin{solution}
                Yes, but the proposition is \FALSE since $x = 6$ doesn't work
            \end{solution}
            
            We can also rewrite this as

            $$\forall x \in \mathbb{R} \mid x + 3 = 10$$
        
    \end{enumerate}
\end{example}

\begin{definition}[Quantifiers]
    The \textbf{existential quantifier} $\exists$ means "there exists"
    \\
    The \textbf{universal quantifier} $\forall$ means "for all" or "for every"
\end{definition}

\begin{example}[Examples of Quantifiers]
    Write each proposition using quantifiers, and determine their truth values
    
    \begin{enumerate}
        \item 
            There exists an integer $x$ such that $x^2 + 2 > 3x$

            \begin{solution}
                $$\exists x \in \mathbb{Z} \mid x^2 + 2 > 3x$$

                \TRUE since $x = 3$ gives $9 + 2 > 9$
            \end{solution}
        

        \item 
            For all real numbers $x$, $x^2 + 2 > 3x$

            \begin{solution}
                $$\forall x \in \mathbb{R} \mid x^2 + 2 > 3x$$

                \FALSE since $x = 1$ gives $1 + 2 > 3$
            \end{solution}

        \begin{note}
            We only need one counter example to disprove a proposition that has $\forall$
        \end{note}

        \item
            There exists an integer $x$ such that $x^3 = 31687$

            \begin{solution}
                $$\exists x \in \mathbb{Z} \mid x^3 = 31687$$            

                \FALSE since $\sqrt[3]{31687} \not \in \mathbb{Z}$
            \end{solution}

        \begin{note}
            We needed to verify that every integer does not satisfy the proposition in order to prove it
        \end{note}
        
    \end{enumerate}
\end{example}

The negations of each quantifier are as follows:

$$\neg \forall \equiv \exists$$
$$\neg \exists \equiv \forall$$

\begin{example}[Negation of Quantifiers]
    Give the negation of the following propositions:

    \begin{enumerate}
        \item
            $\forall x \in \mathbb{R} \mid x^2 + 1 > 0$ \TRUE
            
            \begin{solution}
                $\exists x \in \mathbb{R} \mid x^2 + 1 \leq 0$ \FALSE
            \end{solution}

        \item
            $\exists a \in \mathbb{Z} \mid a^3 + 2 = 3a$ \TRUE $(a = 1)$

            \begin{solution}
                $\forall a \in \mathbb{Z} \mid a^3 + 2 \not = 3a$ \FALSE
            \end{solution}

        \item 
            For every real number $x$, $x^2 + 30 < 11x$ \FALSE $(x=1)$

            \begin{solution}
                $\exists x \in \mathbb{R} \mid x^2 + 30 \geq 11x$ \TRUE
            \end{solution}

        \item 
            For every prime integer $p > 3$, $p^2 - 1$ is divisible by 24 \TRUE

            \begin{solution}
                There exists a prime integer $p > 3$, $p^2 - 1$ is not divisible by 24 \FALSE
            \end{solution}

            \begin{work}
                $$p^2 - 1$$

                \begin{center}
                    We want $p^2 -1$ divisible by both 8 and 3
                \end{center}

                $$(p^2 -1 ) = (p -1)(p + 1)$$
                
                \begin{center}
                    $(p - 1)(p + 1) = (\text{even})(\text{even}$)
                \end{center}

                \begin{itemize}
                    \item p is prime
                    \item $p > 3$
                    \item p is odd
                \end{itemize}
                
            \end{work}
        
    \end{enumerate}
\end{example}

\begin{example}
    \;
    \begin{tablewithheader}
    { |c|c|c| }
    { Statement & When \TRUE & When \FALSE }
    
    $\forall x, P(x)$ & $P(x)$ is \TRUE for all $x$ & There exists at least one $x$ such that $P(x)$ is \FALSE \\
    $\exists x, P(x)$ & There exists at least one $x$ such that $P(x)$ is \TRUE & For all $x$ such that $P(x)$ is \FALSE \\

    \end{tablewithheader}
\end{example}

\begin{example}[Examples with Negation]
    \;
    \begin{tablewithheader}
    { |c|c|c|c| }
    { Statement & Equivalent & When \TRUE & When \FALSE }

    $\neg \forall x, P(x)$ & $\exists x, \neg P(x)$ & There is at least one $x$ such that $P(x)$ is \FALSE & For all $x$, $P(x)$ is \TRUE \\
    $\neg \exists x, P(x)$ & $\forall x, \neg P(x)$ & $P(x)$ is \FALSE for all $x$ & There exists at least one $x$ such that $P(x)$ is \TRUE \\

    \end{tablewithheader}
\end{example}

\begin{example}[Examples of Rewriting with Quantifiers]
    Rewrite the following with Quantifiers

    \begin{enumerate}
        \item 
            For all positive integers, either $6n +1$ or $6n -1$ is prime
            \begin{solution}
                $\forall n \in \mathbb{Z}^+, 6n-1 \text{ is prime } \lor 6n + 1 \text{ is prime }$
                \\
                \FALSE since if $n = 20$ then $6n - 1 = 119$ and $6n + 1 121$ are \textbf{composite} (not prime)
            \end{solution}

        \item 
            There exists an integer $n$ such that $n^2 + n + 41$ is composite
            \begin{solution}
                $\exists n \in \mathbb{Z}, n^2 + n + 41$ is composite
                \\
                \TRUE since if $n = 41$ then 
                
                \begin{align*}
                    n^2 + n + 41 &= 41^2 + 41 + 41 \\
                    &= 41(41 + 1 + 1)
                \end{align*}

                \begin{center}
                    which is composite
                \end{center}
            \end{solution}

    \end{enumerate}
\end{example}

% ========== End the document ==========
\end{document}